\documentclass[11pt]{article}
\usepackage[a4paper,margin=25mm]{geometry}
\usepackage{amsmath,amssymb,amsthm,booktabs,microtype}
\usepackage[hidelinks]{hyperref}
\newtheorem{theorem}{Theorem}
\newtheorem{lemma}[theorem]{Lemma}
\newtheorem{corollary}[theorem]{Corollary}
\newcommand{\SF}{\operatorname{SF}_2}
\newcommand{\diss}{\beta_2}
\newcommand{\NS}{N_S}
\title{A sharp gap bound for two-star-forming sets in bipartite graphs}
\author{}
\date{Working manuscript, 11 November 2026}
\begin{document}
\maketitle
\begin{abstract}
The upper two-star-forming number is the largest cardinality of an
inclusion-minimal set $S$ such that every vertex $x$ outside $S$ belongs to a
not necessarily induced two-leaf star in $G[S\cup\{x\}]$. We prove that its
difference from the dissociation number is at most $\lfloor n/10\rfloor$
in every bipartite graph of order $n$. The coefficient is sharp even for
two-connected graphs of minimum degree two and maximum degree four. The
proof constructs two dissociation sets from disjoint deletion witnesses.
We also prove that both parameters are additive over copies of a ten-vertex
graph under permitted bipartite interconnections through four specified vertices.
For graphs with gap $n/10>0$, we determine the intrinsic layers and
the unique maximum-cardinality minimal forming set, and give an exact
neighbourhood recognition criterion. For complete cores with at least
four vertices in each part this reduces to paired witness-neighbourhood
inequalities; the ten-vertex case is classified separately. We also prove that the asymptotic
coefficient remains sharp under every fixed minimum-degree constraint.
\end{abstract}

\section{Introduction}

How far can the upper two-star-forming number of a bipartite graph exceed
its dissociation number? We prove that the difference is at most one tenth
of the order, and construct two-connected graphs of maximum degree four
attaining this bound at arbitrarily large orders. Thus the restriction to
bipartite graphs gives a sharp quantitative comparison between the two
parameters, although they need not be equal.

All graphs are finite and simple. For a bipartition $(A,B)$ write
$Y_A=Y\cap A$, $Y_B=Y\cap B$, and
$N(T)=\bigcup_{v\in T}N_G(v)$. A set $S\subseteq V(G)$ is
\emph{two-star-forming} if every $x\notin S$ lies in a copy of $K_{1,2}$
in $G[S\cup\{x\}]$; the copy need not be induced. Let $\SF(G)$ denote
the largest cardinality of an inclusion-minimal such set. A
\emph{dissociation set} induces maximum degree at most one, and $\diss(G)$
is its maximum cardinality. Chellali and Favaron~\cite{CF09} established
the comparison $\diss(G)\le\SF(G)$ in the more general framework of
$k$-star-forming sets. The parameters measure different forms of maximality:
a dissociation set restricts its own induced subgraph, whereas minimal
star formation requires each single deletion to destroy formation, possibly
because the deleted vertex itself becomes uncovered.

Equality for bipartite graphs, and for triangle-free graphs at $k=2$, was
proposed in Conjectures 4.1, 4.2 and 7.2 of Sahbi's version~2
preprint~\cite{Sahbi}; that preprint was subsequently withdrawn. The
ten-vertex graph in Section~\ref{sec:attainment} has $\diss=5$ and
$\SF=6$. Our main result determines the largest possible difference per
vertex, rather than an equality of the two parameters.

Favaron~\cite[Theorem~10]{Favaron12} gives
$SF_k(G)\le n-\delta+k-1$ for $1\le k\le\delta$, where $\delta$ is the
minimum degree. That order bound does not compare $\SF$ with $\diss$.
Bhattacharjee~\cite{Bhattacharjee} gives higher-$k$ bipartite examples,
chain-graph equalities and complexity results for the threshold decision
problem $\SF(G)\ge t$. Attachment
constructions preserving extremal dissociation numbers also appear in
Bock et al.~\cite{BPPP} and Li~\cite{Li}. Here the global inequality comes
from two exchanges applied to every minimal forming set; the attaining
construction controls both parameters under its permitted interfaces.

\begin{theorem}\label{thm:gap}
Every bipartite graph $G$ of order $n$ satisfies
\[
  \SF(G)-\diss(G)\le\left\lfloor\frac n{10}\right\rfloor.
\]
The coefficient $1/10$ is sharp among two-connected bipartite graphs of
minimum degree two and maximum degree four. If
$\SF(G)-\diss(G)=n/10>0$, then exactly $n/5$ vertices have degree two.
\end{theorem}

The proof of the upper bound and equality condition occupies
Section~\ref{sec:bound}. Section~\ref{sec:attainment} supplies the exact
composition theorem and the attaining graphs. Section~\ref{sec:equality}
determines the intrinsic equality structure and gives equality graphs
without a spanning $H$-factor. Section~\ref{sec:degree} proves degree
refinements and minimum-degree asymptotic sharpness.

\section{Two exchanges from deletion witnesses}\label{sec:bound}

Write $\NS(x)=N_G(x)\cap S$ and $d_S(x)=|\NS(x)|$. The definition of
two-star formation is equivalent to
\begin{equation}\label{eq:local}
 x\notin S\quad\Longrightarrow\quad
 d_S(x)\ge2\ \text{ or }\
 \exists y\in\NS(x):d_S(y)\ge1.
\end{equation}
The alternatives make $x$ the centre or a leaf of the star. Formation is
upward closed: enlarging a forming set preserves a witnessing star at
each vertex still outside. Consequently, inclusion-minimality is
equivalent to failure of formation after every single deletion.

Fix an inclusion-minimal forming set $S$ and a bipartition $(A,B)$ of $G$.
Let $R$ consist of the vertices in components of $G[S]$ of order one or two,
let $C=S\setminus R$, and let $I$ be the set of isolated vertices of $G[S]$.
Every $v\in C$ lies on a three-vertex path in $G[S]$. After deleting $v$,
the newly outside vertex $v$ therefore satisfies~\eqref{eq:local}.
Minimality supplies a vertex $x\notin S$ that loses coverage. We use the
local certificate restrictions of~\cite[Section~5]{Sahbi}, giving their
short proof here. There are three possible forms:
\begin{align}
 &\NS(x)=\{v\};\label{eq:type1}\\
 &\NS(x)=\{v,i\},\qquad i\in I;\label{eq:type2}\\
 &\NS(x)=\{\ell\},\qquad \NS(\ell)=\{v\}.\label{eq:type3}
\end{align}
Indeed, if $xv$ is an edge, $x$ retains at most one selected neighbour.
If this neighbour $i$ exists, it is isolated after the deletion. Since
$i$ and $v$ lie in the same bipartition class, they are nonadjacent,
so $i$ was isolated already. If $xv$ is not an edge, the selected
neighbourhood of $x$ is unchanged. It can lose coverage only through a
unique neighbour $\ell$ whose sole selected neighbour was $v$.

For every $c\in C$ admitting~\eqref{eq:type3}, choose one triple
$(c,\ell_c,x_c)$, and let $p$ count the chosen support vertices. Each support has
selected degree at least two; otherwise it and $\ell_c$ form a component
of order two. The chosen leaves have distinct unique selected neighbours,
and cannot themselves be supports. The outside vertices $x_c$ are distinct
because their unique selected neighbours are distinct. Thus these triples
are vertex-disjoint.

Let
\[
 U=C\setminus\bigl(\{c:\text{chosen support}\}
                 \cup\{\ell_c:\text{chosen leaf}\}\bigr),
 \qquad u=|U|.
\]
For every $v\in U$, choose a witness $x_v$ of
form~\eqref{eq:type1} or~\eqref{eq:type2}. No vertex of $U$ admits
form~\eqref{eq:type3}, since all such vertices were chosen as supports.
Each $x_v$ has exactly one neighbour in $C$, namely $v$. These witnesses
are distinct from one another and from the triple witnesses. We have
\begin{equation}\label{eq:partition}
 |S|=|R|+2p+u.
\end{equation}

For $i\in I$, let $h(i)$ count the chosen witnesses $x_v$ adjacent to $i$.
Let $a$ count the isolates with $h(i)=1$ and $b$ those with $h(i)\ge2$.
Each $x_v$ meets at most one isolate, so $u\ge a+2b$. The $p$ triples,
the $u$ pairs $(v,x_v)$ and the $a+b$ touched isolates are mutually
disjoint. It follows that
\begin{equation}\label{eq:budget}
 n\ge3p+2u+a+b\ge3p+3a+5b.
\end{equation}

We construct two dissociation sets, then repeat each construction with
$A$ and $B$ exchanged. In the first, retain $R$ except the touched
isolates in $A$. From each triple select $c,\ell_c$ when $c\in A$,
and $\ell_c,x_c$ when $c\in B$. From each remaining pair $(v,x_v)$,
select both vertices when $v\in A$ and neither when $v\in B$.
Denote the resulting set by $D_A$.

All selected outside witnesses lie in $B$, so none are adjacent to one
another. A selected triple witness sees only its chosen leaf in $S$.
A selected pair witness sees its owner $v$, and its optional isolate
neighbour lies in $A$ and was removed. Among vertices of $C\cap B$,
only chosen leaves with supports in $A$ remain. Such a leaf has its support
as its unique selected neighbour in $S$. There are no selected edges
between $R$ and $C$, and no chosen witness meets an order-two component
of $G[S]$. These observations show that $G[D_A]$ has maximum degree at most one.
Counting both orientations gives
\begin{equation}\label{eq:first}
 \diss(G)\ge\frac{|D_A|+|D_B|}{2}=|S|-\frac{a+b}{2}.
\end{equation}

For the second construction, let $T$ be $S$ together with all $p+u$
chosen outside witnesses. Select $T\cap A$, together with $R\cap B$
except its isolates with $h(i)\ge2$, and call this set $E_A$.
A retained endpoint of an order-two component in $R\cap B$ sees only
its mate; no chosen witness meets such a component. A retained isolate
in $R\cap B$ meets at most one chosen witness. On the $A$ side, a vertex
of $S$ sees at most one vertex of $R$, and a chosen witness sees at most
one isolate of $R$. Hence $E_A$ is a dissociation set. Since
$|T|=|S|+p+u$, equation~\eqref{eq:partition} gives
\begin{equation}\label{eq:second}
 \diss(G)\ge\frac{|E_A|+|E_B|}{2}
 =\frac{|T|+|R|-b}{2}=|S|-\frac{p+b}{2}.
\end{equation}
These are alternative constructions; their resources are counted
separately, while~\eqref{eq:budget} counts each vertex once.

Combining~\eqref{eq:budget}, \eqref{eq:first} and~\eqref{eq:second},
\begin{equation}\label{eq:chain}
 10\bigl(|S|-\diss(G)\bigr)
 \le5\min(a,p)+5b\le3a+3p+5b\le n.
\end{equation}
The middle inequality follows from
$3a+3p-5\min(a,p)=\min(a,p)+3|a-p|\ge0$.
Taking $|S|=\SF(G)$ and using integrality proves the upper bound.

Suppose now that $\SF(G)-\diss(G)=n/10>0$. All inequalities in
\eqref{eq:chain} and~\eqref{eq:budget} are equalities. Equality in
$5\min(a,p)\le3(a+p)$ forces $a=p=0$. Hence $u=2b$ and $n=5b$.
The counted vertices exhaust $V(G)$: all outside vertices are pair
witnesses, and $R$ consists exactly of the $b$ touched isolates.
Their witness incidences total at most $u=2b$ and each isolate has at
least two, so every one has exactly two neighbours. They have no selected
neighbours and there are no other outside vertices. These $b=n/5$
vertices have degree two. Theorem~\ref{thm:intrinsic} below proves
that these are all degree-two vertices. The assertion concerns equality with $n/10$,
rather than equality with its integer part.

\section{Safe deletion, exact composition and attainment}\label{sec:attainment}

A sufficient condition for safe composition can be stated for an arbitrary
induced block $K$ and a boundary $P\subseteq V(K)$ containing every vertex
of $K$ with a neighbour outside $K$.

\begin{lemma}\label{lem:boundary}
Let $S$ form an ambient graph and put $C=S\cap V(K)$. A vertex
$v\in C\setminus P$ can be deleted safely if $C\setminus\{v\}$ forms
$K$ and every selected boundary vertex having a neighbour in $C$ retains
one in $C\setminus\{v\}$.
\end{lemma}
\begin{proof}
All vertices of $K$ outside the new selected set are covered
internally. A vertex $x$ beyond $K$ is not adjacent to $v$. Its old star
can be lost only when $v$ was a selected leaf, the selected centre was
$p\in P$, and $x$ was the other leaf. A retained internal selected
neighbour of $p$ replaces $v$ in that star.
\end{proof}

Here is a finite sufficient test for a cap $s$ on $|S\cap V(K)|$.
For every $C\subseteq V(K)$ of size greater than $s$, impose the local
precondition
\begin{equation}\label{eq:boundarypre}
 \forall x\in V(K)\setminus(C\cup P):\quad
 |N_K(x)\cap C|\ge2\quad\text{or}\quad
 \exists y\in N_K(x)\cap C:
       (N_K(y)\cap C\ne\varnothing\ \text{or }y\in P).
\end{equation}
Every global forming set satisfies this condition: $x$ has no external
neighbour, and its unique selected neighbour can obtain external activity
only through $P$. If every such $C$ admits a deletion as in
Lemma~\ref{lem:boundary}, no global minimal forming set can use more
than $s$ vertices of $K$. This proves a cap independent of the ambient
graph, without requiring outside activity to be internal.

\begin{theorem}\label{thm:generalcomposition}
Start with finitely many pairwise vertex-disjoint graphs $K_i$ and
boundaries $P_i\subseteq V(K_i)$. Form $G$ on
$V(G)=\bigsqcup_i V(K_i)$ by retaining all block edges and adding only
edges between $P_i,P_j$ for distinct $i,j$. Suppose their global caps are $s_i$.
For each $i$, suppose $C_i$ forms $K_i$, $|C_i|=s_i$, and for every
$v\in C_i$ there is a vertex $x\in V(K_i)\setminus P_i$, outside $C_i\setminus\{v\}$,
whose neighbourhood in $C_i\setminus\{v\}$ is either empty or is
$\{y\}$, where $y\notin P_i$ has no selected neighbour after deletion.
Then $\SF(G)=\sum_i s_i$.
If each block has bipartition sides of size $b_i$, $\diss(K_i)=b_i$,
and interfaces respect these sides, then also $\diss(G)=\sum_i b_i$.
\end{theorem}
\begin{proof}
The union of the $C_i$ forms $G$ by internal stars. Its deletion
witnesses remain failures: $x$ gains no external neighbours, and in the
one-neighbour case neither does $y$. It is therefore minimal. The block
caps give the opposite inequality. For dissociation, restriction to each
block gives the upper bound and a whole global side gives the lower bound.
\end{proof}

The boundary safeguard in Lemma~\ref{lem:boundary} has content. On the
path $p$--$v$--$a$--$b$, take $C$ to be all four vertices and $P=\{p\}$.
After deleting $v$, the remaining set still forms the path. Attach one
outside vertex $x$ adjacent only to $p$. Its former star $x$--$p$--$v$
is lost, and $x$ becomes uncovered. Internal formation alone does not
make the deletion globally safe.


Let $H$ have parts $A_H=\{a_1,\ldots,a_5\}$,
$B_H=\{b_1,\ldots,b_5\}$ and cross-edge matrix
\begin{equation}\label{eq:matrix}
 \begin{pmatrix}
 1&1&0&0&1\\1&1&0&1&0\\0&0&0&1&1\\
 0&1&1&1&1\\1&0&1&1&1
 \end{pmatrix}.
\end{equation}
The attachment vertices are $P=\{a_1,a_2,b_1,b_2\}$.

\begin{theorem}\label{thm:composition}
Take $r$ disjoint induced copies of $H$, and add any collection of edges
between distinct copies, always joining a vertex labelled $b_1$ or $b_2$
to one labelled $a_1$ or $a_2$. The resulting graph $F$ satisfies
\[
 \diss(F)=5r,\qquad \SF(F)=6r.
\]
\end{theorem}

We prove a local upper bound that also applies in a nonbipartite
surrounding graph. First, if a selection $T$ in $H$ has at least six
vertices, formation is equivalent to domination. To see the converse,
suppose an outside vertex $x$ has a unique selected neighbour $y$ and
$y$ has no selected neighbour. The disjoint sets $N_H(x)\setminus\{y\}$
and $N_H(y)$ lie in $V(H)\setminus T$. Every edge of $H$ has degree sum
at least six, so their union has at least five vertices, a contradiction.
Disjointness follows from bipartiteness; the degree-two vertices have
only degree-four neighbours and all other degrees are at least three.

The closed neighbourhoods of size at most four are
\begin{align*}
 &\{a_3,b_4,b_5\},\quad\{b_3,a_4,a_5\},\\
 &\{a_1,b_1,b_2,b_5\},\quad\{a_2,b_1,b_2,b_4\},\\
 &\{b_1,a_1,a_2,a_5\},\quad\{b_2,a_1,a_2,a_4\}.
\end{align*}
The two triples are disjoint. Any two distinct displayed quadruples
intersect in at most two vertices. All other closed neighbourhoods have
size five.

\begin{lemma}\label{lem:localcap}
Suppose $H$ is induced in a graph $G$ and all edges from $H$ to the rest
of $G$ meet $H$ at $P$. Every inclusion-minimal forming set $S$ satisfies
$|S\cap V(H)|\le6$.
\end{lemma}
\begin{proof}
Suppose $C=S\cap V(H)$ has at least seven vertices and put
$W=V(H)\setminus C$, so $|W|\le3$. The set $C$ dominates $H$.
Otherwise some $N_H[x]$ lies in $W$; only $a_3,b_3$ can have such small
closed neighbourhoods. Neither is an attachment vertex, so $x$ has no
selected neighbour in $G$, contradicting formation.

At most two candidates for deletion destroy domination. Completing one
of the disjoint triples from $W$ requires two elements of $W$, so at
most one candidate completes a triple. Completing a quadruple requires
$|W|=3$ and $W$ contained in that quadruple. The intersection bound
permits at most one such quadruple. A size-five neighbourhood is too large.

Let $t=|C\cap P|$. If $t\le2$, there are at least $7-t$ selected
nonattachment candidates. Besides the two domination failures, at most
$t$ candidates destroy the existing internal selected-neighbour activity
of an attachment: each selected attachment has at most one unique selected
neighbour. Since $7-t>2+t$, some candidate is safe. If $t\ge3$, the
selected attachment vertices induce a complete bipartite graph with both
sides nonempty. Every selected attachment therefore has a selected
attachment neighbour, which no nonattachment deletion removes. There
are at least three candidates and at most two domination failures, so
again a safe candidate $v$ exists.

The set $C\setminus\{v\}$ dominates $H$ and has at least six vertices,
so it is forming inside $H$. Each selected attachment with an internal
selected neighbour before deletion retains one. Every vertex of
$V(H)\setminus(C\setminus\{v\})$, including $v$, therefore has a
witnessing star entirely inside $H$ after deletion. For an
outside vertex $x$ of $G$ beyond the copy, an old star avoiding $v$
survives. A star using $v$ has $v$ as a leaf, because $xv$ is not an
edge. Its selected centre is an attachment, and its other leaf is $x$.
Replace $v$ by a surviving selected internal neighbour of that centre.
This proves formation of $S\setminus\{v\}$, contradicting minimality.
\end{proof}

\begin{proof}[Proof of Theorem~\ref{thm:composition}]
Lemma~\ref{lem:localcap} gives $\SF(F)\le6r$. In each copy select
$S_0=\{a_1,a_2,a_3,b_1,b_2,b_3\}$. Every omitted vertex has two selected
neighbours. Single deletions fail at the following vertices:
\begin{center}
\begin{tabular}{@{}lcccccc@{}}
\toprule
Deleted & $a_1$ & $a_2$ & $a_3$ & $b_1$ & $b_2$ & $b_3$\\
Failure & $b_5$ & $b_4$ & $a_3$ & $a_5$ & $a_4$ & $b_3$\\
Remaining neighbours & $a_3$ & $a_3$ & none & $b_3$ & $b_3$ & none\\
\bottomrule
\end{tabular}
\end{center}
All these failure vertices are nonattachments, and $a_3,b_3$ remain
isolated in the selected subgraph under every interface. Thus the union
of the $S_0$ is minimal, proving $\SF(F)=6r$.

The whole $A_H$ side is independent in $F$, giving $\diss(F)\ge5r$.
It remains to exclude a dissociation six-set in $H$. By the symmetry
of the matrix under exchanging the sides, its side split is $(1,5)$,
$(2,4)$ or $(3,3)$. The first fails minimum degree two. In the second,
omitting only one $B_H$ vertex permits selected degree at most one only
at $a_3$, so two selected $A_H$ vertices are impossible. In the third,
the degree-four vertices $a_4,a_5$ cannot be selected; the $A_H$ side is
therefore $\{a_1,a_2,a_3\}$. Both $a_1,a_2$ must avoid at least two of
their three neighbours. With only two omitted $B_H$ vertices this forces
both $b_1,b_2$ omitted. The remaining $b_4,b_5$ both meet $a_3$, a
contradiction. Every global dissociation set consequently uses at most
five vertices of each copy, proving $\diss(F)=5r$.
\end{proof}

\begin{corollary}\label{cor:topology}
Let $Q$ be a connected simple graph with no bridge, with an orientation
having at most two incoming and two outgoing edges at each vertex.
There is a two-connected bipartite graph $F$ of order $10|V(Q)|$,
maximum degree four and minimum degree two, with
$\SF(F)-\diss(F)=|V(F)|/10$.
\end{corollary}
\begin{proof}
Replace each vertex $i$ of $Q$ by $H_i$. Assign distinct outgoing edges
to $b_1^{(i)},b_2^{(i)}$ and distinct incoming edges to
$a_1^{(i)},a_2^{(i)}$, then join the assigned vertices along each edge.
Each attachment gains at most one edge, so all degrees lie between two
and four. Each copy has the Hamilton cycle
\[
 a_1,b_1,a_2,b_2,a_4,b_3,a_5,b_4,a_3,b_5,a_1.
\]
Deleting a vertex leaves its copy connected and removes at most one
interface edge. Since $Q$ has no bridge, the surviving interfaces still
connect the copies. The graph remains connected after every vertex
deletion, proving two-connectivity. Theorem~\ref{thm:composition} gives
the exact gap.
\end{proof}

Taking $Q$ to be a cycle gives a cyclic attaining family. More generally,
every connected simple four-regular $Q$ qualifies: an Euler orientation
has two edges entering and leaving each vertex, and a connected graph
with all degrees even has no bridge. For example $Q=C_m^2$, $m\ge5$,
gives attaining graphs of order $10m$ and size $18m$. Here cyclic
distances one and two define the square; $C_5^2=K_5$. All four attachments
of every copy then have external neighbours. In these examples the two
degree-two vertices per copy meet the necessary count in
Theorem~\ref{thm:gap} exactly.

\section{Intrinsic structure of equality}\label{sec:equality}

Throughout this section $G$ is bipartite, and equality means $\SF-\diss=n/10>0$, rather than
attainment of the integer part.

\begin{theorem}\label{thm:intrinsic}
Every equality graph has order $10k$, $\SF=6k$ and $\diss=5k$.
Let $I$ be its set of degree-two vertices, $X=N_G(I)$ and
$C=V(G)\setminus(I\cup X)$. These are disjoint sets. The neighbourhoods
of vertices of $I$ are disjoint pairs partitioning $X$, and the edges
between $C$ and $X$ form a perfect matching. The only other edges are
inside $C$ and inside $X$. Each bipartition side has $k$ vertices in
$I$ and $2k$ in each of $C,X$. Moreover,
$\delta(G[C])\ge2$ and $\delta(G[X])\ge1$.
The unique maximum-cardinality minimal forming set is $I\cup C$.
In particular, exactly $n/5$ vertices have degree two.
\end{theorem}
\begin{proof}
Apply the equality ledger at the end of Section~\ref{sec:bound} to any
maximizing minimal forming set $S$. Initially write $I=R$, $C=S\setminus I$
and $X=V\setminus S$. Every vertex $c\in C$ has a distinct witness
$x_c\in X$ with selected neighbourhood $\{c,i(c)\}$, and every vertex
of $I$ meets exactly two witnesses. All outsiders are these witnesses.
Thus the claimed pair and matching incidences hold. We have
$|I|=b$, $|C|=|X|=2b$, $|S|=3b$ and $\diss=5b/2$, so $b=2k$.
If $I$ has $r,s$ vertices in the two sides, the side sizes are
$3r+2s$ and $2r+3s$. Each side is independent of size at most $5k$,
so $r=s=k$.

Suppose $c\in C_A$ has at most one neighbour in $C_B$. Start with the
entire side $B$, remove $x_c$, and add $c$ and $i(c)$. The vertex $c$
sees at most its one core neighbour, while $i(c)$ sees only its other
witness. These possible partners are distinct. The result is a
dissociation set of size $5k+1$, a contradiction. Thus every core
vertex has core degree at least two and full degree at least three.

If $x\in X_A$ has no neighbour in $X_B$, the set
$C_A\cup X_B\cup I_B$ of size $5k$ induces just the matching
$C_A$--$X_B$. Adding $x$ pairs it only with its unique neighbour in
$I_B$, giving the same contradiction. Thus each witness has full
degree at least three. The originally selected isolates are therefore
exactly the degree-two vertices of $G$. Their neighbours determine
$X$, and the remaining vertices determine $C$, independently of $S$.
Repeating the argument for any maximizing minimal set proves uniqueness.
\end{proof}

We next give an exact adjacency criterion. For $T\subseteq A$, let
$p(T)$ count vertices $v\in T$ having a private neighbour, that is, a
vertex $w\in B$ with $N(w)\cap T=\{v\}$. Count each owner $v$ once,
even when it has several private neighbours. Fixing the $A$ part of a
dissociation set gives
\begin{equation}\label{eq:private}
 \diss(G)=|B|+\max_{T\subseteq A}\bigl(|T|+p(T)-|N(T)|\bigr).
\end{equation}
Indeed one may include all $B$ vertices with no neighbour in $T$, none
with two or more, and one private neighbour for each owner admitting one.
Those choices are mutually compatible because $B$ is independent.

\begin{theorem}\label{thm:recognition}
A nonempty bipartite graph has positive unrounded equality if and only
if the following conditions hold. Its degree-two set $I$ is nonempty
and independent; each vertex of $X=N(I)$ has exactly one neighbour in
$I$; the $C$--$X$ edges form a perfect matching, where
$C=V\setminus(I\cup X)$; and there is a balanced bipartition $(A,B)$
satisfying
\[
 |N(T)|\ge |T|+p(T)\qquad(T\subseteq A).
\]
\end{theorem}
\begin{proof}
Necessity follows from Theorem~\ref{thm:intrinsic} and~\eqref{eq:private}.
Conversely put $b=|I|$. The pair and matching conditions give
$|X|=|C|=2b$ and $n=5b$. The set $S=I\cup C$ is forming: every
omitted vertex has its isolate and core owner as selected neighbours.
Deleting an isolate leaves that newly outside vertex with no selected
neighbour. Deleting a core owner leaves its witness with only its
isolated selected neighbour, so that witness is neither a star centre
nor a covered leaf. Thus $S$ is minimal of size $3b$.
Equation~\eqref{eq:private} and the balanced bipartition give
$\diss=5b/2$. The universal bound forces $\SF\le3b$, proving equality.
Integrality forces $b$ even.
\end{proof}

This is an exact all-subsets test, rather than a claim of polynomial-time
recognition. Fix $k\ge1$. We define the following template independently of equality.
Let $G$ be bipartite with prescribed layers $V(G)=I\sqcup C\sqcup X$,
where each side contains $k$ vertices of $I$ and $2k$ vertices of each
of $C,X$. The edges incident with $I$ are exactly disjoint two-leaf
stars centred at $I$, whose leaf pairs partition $X$. The $C$--$X$
edges form a perfect matching. All remaining edges are arbitrary
edges within $G[C]$ and within $G[X]$. Initially $I$ need not be the
entire degree-two set, and no minimum degree is imposed on the two
internal graphs. A shorter criterion controls all dissociation sets
avoiding one core side in this template.
Let $\mathcal P_A,\mathcal P_B$ be the partitions of $X_A,X_B$ into
pairs centred at vertices of $I_B,I_A$, respectively. For
$\mathcal J\subseteq\mathcal P_B$, define
\[
 z(\mathcal J)=\bigl|\{P\in\mathcal P_A:
 P\subseteq N_{G[X]}(\bigcup\mathcal J)\}\bigr|,
\]
and define the reversed quantity by exchanging the sides.

\begin{theorem}\label{thm:paired}
For every such template, with arbitrary core edges,
\[
 \max_{D\text{ dissociation},\ D\cap C_B=\varnothing}|D|
 =5k+\max_{\mathcal J\subseteq\mathcal P_B}
          (|\mathcal J|-z(\mathcal J)).
\]
Consequently $z(\mathcal J)\ge|\mathcal J|$ in both orientations is
necessary for equality. If $k\ge2$ and $G[C]=K_{2k,2k}$, these two
conditions are also sufficient.
\end{theorem}
\begin{proof}
A dissociation set avoiding $C_B$ can be enlarged or exchanged to
contain all $C_A$. For a missing owner $c$, add it if its witness is
absent or isolated; otherwise replace that witness by $c$. The owner
has no other available neighbour, so size never decreases.

Now put $T=D\cap X_B$. Each selected witness is paired with its core
owner, excluding its centre in $I_A$ and all its neighbours in $X_A$.
Let $t(T)$ count the witness pairs touched by $T$, $f(T)$ those wholly
selected, and $z(T)$ the opposite pairs wholly contained in $N_{G[X]}(T)$.
The available $I_A$ centres contribute $k-t(T)$. Each opposite centred
pair contributes two vertices if a witness is available, and one if
both are forbidden. These choices are compatible and attain exactly
\[
 2k+|T|+k-t(T)+2k-z(T)=5k+f(T)-z(T).
\]
Removing singleton selections from $T$ preserves $f(T)$ and cannot
increase $z(T)$, so a maximum is attained by a union of whole pairs.
This proves the formula and its reversed version.

In a complete core, a dissociation set meeting both core sides uses
at most one vertex of each. Each of the $2k$ disjoint $I$-centred
three-vertex paths contributes at most two vertices, giving
$|D|\le2+4k\le5k$. The paired conditions handle the other cases by
the formula. A whole side attains $5k$, and the designated minimal
forming set $I\cup C$ and the universal bound give $\SF=6k$.
\end{proof}

Bhattacharjee~\cite[final computational remark]{Bhattacharjee} reports by exhaustive
search that $H$ is the unique bipartite counterexample through order ten.
The following gives a structural proof of its order-ten uniqueness.

\begin{corollary}
Up to isomorphism, $H$ is the unique ten-vertex bipartite graph with
$\SF-\diss=1$.
\end{corollary}
\begin{proof}
The intrinsic layers have one isolate and two owners and witnesses on
each side. Relabel their fixed incidences as in~\eqref{eq:matrix}.
All other edges lie in the core or witness subgraph, so the graph is a
spanning subgraph of $H$. The group generated by the automorphisms
$(a_1\ a_2)(b_4\ b_5)$ and $(b_1\ b_2)(a_4\ a_5)$ act transitively
on each of its four core edges and four witness edges.
Deleting $a_1b_1$ admits the dissociation set
$\{a_1,a_2,a_3,a_4,b_1,b_3\}$, whose edges are $a_2b_1,a_4b_3$.
Deleting $a_4b_4$ admits $\{a_1,a_3,a_4,b_1,b_3,b_4\}$, whose edges
are $a_1b_1,a_3b_4,a_4b_3$. Further deletions preserve these sets.
Any missing edge thus contradicts $\diss=5$, proving the assertion.
\end{proof}

An $H$-factor means a spanning collection of vertex-disjoint, not
necessarily induced copies of $H$. There are equality graphs with no
$H$-factor.
For $r\ge2$, start with a core $K_{2r,2r}$ on vertices
$a_{j,t},b_{j,t}$, $1\le j\le r$, $t\in\{0,1\}$.
For each $j$ add the six-cycle
\[
 i_j,y_{j,0},x_{j,0},k_j,x_{j,1},y_{j,1},i_j
\]
and the owner edges $a_{j,t}y_{j,t}$ and $b_{j,t}x_{j,t}$. There are
no other edges. Denote this graph by $L_r$.

The paired conditions alone do not handle arbitrary cores. Take $k$
disjoint $K_{2,2}$ cores, match each witness pair to its aligned opposite
pair, and add the prescribed owner and isolate edges. Both conditions
hold with equality. In each ten-vertex module, with the notation just defined, $\{i_j,k_j,y_{j,0},x_{j,1},a_{j,1},b_{j,0}\}$ induces three
disjoint edges. Hence $\diss\ge6k>5k$.

\begin{corollary}
For every $r\ge2$, $L_r$ is two-connected, has order $10r$,
$\diss(L_r)=5r$ and $\SF(L_r)=6r$, and has no spanning $H$-factor.
\end{corollary}
\begin{proof}
The witness interface is a pair-preserving perfect matching, so
$z(\mathcal J)=|\mathcal J|$ in both orientations.
Theorem~\ref{thm:paired} gives the exact parameters. The core remains
connected after one deletion; each six-cycle remains connected after
an internal deletion and retains at least three owner attachments.
A core deletion also leaves three attachments at its cycle. Thus every
single vertex deletion leaves the whole graph connected.

Degrees are two at $i_j,k_j$, three at $x_{j,t},y_{j,t}$ and $2r+1$
at core vertices. An $H$ copy containing a degree-two vertex of $L_r$
would have to use it as a degree-two vertex of $H$. Both its neighbours
in $H$ have degree four, while both actual neighbours in $L_r$ have
degree three. This is impossible even for a non-induced copy. Such
vertices cannot be covered by an $H$-factor.
\end{proof}

\section{Degree constraints and asymptotic sharpness}\label{sec:degree}

\begin{theorem}\label{thm:degree}
Let $G$ be bipartite, let $g=\SF(G)-\diss(G)$, and let $n_2$ count degree-two vertices.
If $2\le\delta(G)\le\Delta(G)\le D$, then
\[
 (10D+2)g\le Dn+n_2.
\]
If in addition $\delta(G)\ge d\ge2$, then
\[
 (10D+2d-4)g\le Dn.
\]
\end{theorem}
\begin{proof}
Use the witness ledger for a maximizing minimal set, and put
\[
 m=\min(a,p),\quad e=u-a-2b,\quad r=|R|-a-b,\quad
 q=n-|S|-p-u,\quad t=m+b-2g.
\]
These quantities are nonnegative. Exact counting gives
\[
 L:=n-10g=r+(3a+3p-5m)+2e+q+5t\ge2e+q.
\]
Touched isolates have no selected neighbours. Triple witnesses meet
none of them, each pair witness meets at most one, and each of the $q$
remaining outsiders supplies at most $D$ incidences. Therefore
\[
 \sum_{i\text{ touched}}\deg_G(i)\le u+Dq.
\]
At most $n_2$ of the $b$ multiply touched isolates have degree two.
The left side is at least $2(a+b)+(b-n_2)$, yielding
$a+b-n_2\le e+Dq$. Since $2g\le m+b\le a+b$, it follows that
$2g\le n_2+e+Dq\le n_2+DL$, proving the first assertion.

Under minimum degree $d$, the incidence bound instead gives
$(d-1)a+(d-2)b\le e+Dq$. Since $m\le a$,
$2(d-2)g\le(d-2)(m+b)\le e+Dq\le DL$,
which proves the second assertion.
\end{proof}

These estimates are refinements of the universal bound, without a claim
of sharpness for every fixed degree range. A fixed minimum degree alone
has a different answer from exact equality.

\begin{theorem}\label{thm:mindegree}
For every fixed integer $d\ge3$, there are two-connected bipartite graphs with
minimum degree exactly $d$ and
$(\SF-\diss)/n\longrightarrow1/10$.
\end{theorem}
\begin{proof}
Put $s=d-2$ and take $r\ge d$ copies of $H$. In copy $j$, write
$I_j=\{a_3,b_3\}$, $C_j=\{a_1,a_2,b_1,b_2\}$ and
$X_j=\{a_4,a_5,b_4,b_5\}$. Complete all cross-edges between the core
sides over all copies, and separately between the witness sides.
Add $s$ new vertices on each side, $Q_A,Q_B$, joining every $Q_A$
vertex to every $b_3$ and every $Q_B$ vertex to every $a_3$.
The induced copies of $H$ are preserved. Degrees at $I,C,X,Q$ are
$d,2r+1,2r+2,r$, respectively; the order is $n=10r+2s$.

The union of the original six-sets is minimal forming of size $6r$.
Witnesses retain their selected neighbourhoods, and each new $Q$ vertex
has $r$ selected neighbours. Deleting an owner leaves its witness with
only an isolated selected neighbour; deleting an isolate leaves that
newly outside vertex with no selected neighbour. The added edges do
not change either failure.

We claim $\diss=5r+s$. A dissociation set using at most one $Q$ class
has at most five vertices in each induced $H$ and at most $s$ new
vertices. If it uses both classes, at most one $a_3$ and one $b_3$ can
be selected globally. Thus at most two copies contain a selected
isolate. In any other copy, the two disjoint $K_{2,2}$ subgraphs on
$C_j,X_j$ permit at most four selected vertices. The total is at most
$4r+2+2s\le5r+s$, since $r\ge s+2$. A whole side attains $5r+s$.

Each $H$ has the Hamilton cycle used in Corollary~\ref{cor:topology},
so deletion of a block vertex leaves its block connected and removes
at most one of its two kinds of $Q$ attachment. If a $Q$ vertex is
deleted and $s=1$, the other $Q$ class still connects all blocks; if
$s>1$, both classes remain. All surviving $Q$ vertices have neighbours
in at least $r-1$ blocks. This proves two-connectivity.
Finally $g\ge r-s$ and $g\le n/10$, giving the asserted limit.
\end{proof}


\appendix
\section{One budget for overlapping Hall-tight regions}

We record separately a simultaneous counting lemma for the grid
compensation problem. Its hypotheses refer to terminal alternating paths;
its conclusion does not assert that an endpoint assignment can be routed.
Let an exterior bipartite graph have white side $T$, black side $B$,
$|T|=|B|+1$, and a matching saturating $B$. There are $r$ additional
black centres $C$ and $r-1$ additional white sources $W$.
Suppose a full auxiliary matching, together with an actual matching $M$,
decomposes into alternating terminal paths defining a bijection
\[
 \pi:T\sqcup W\longrightarrow B\sqcup C.
\]
A path starting in $T$ is called TH, and one starting in $W$ is HH.
Let $k$ count the TH paths ending at centres. A direct TH path is an
edge of the exterior graph; centres have no direct $T$ neighbour.
For the expense conclusion assume actual-edge weights are at least one,
every non-direct TH path contains an actual edge, and actual edges on
distinct terminal paths are disjoint. Let $f$ be the sum of actual-edge
weights on all TH paths. These hypotheses hold in the separated-core
setting of the compensation problem, and are stated here explicitly.

Choose finitely many nonempty tight sets $D_j\subseteq B$ with
$F_j=N_T(D_j)$ and $|F_j|=|D_j|$, allowing overlap and nesting.
For nonnegative weights $\lambda_j$, not all zero, put
\[
 \phi_B(b)=\sum_j\lambda_j\mathbf1_{b\in D_j},\qquad
 \phi_T(t)=\sum_j\lambda_j\mathbf1_{t\in F_j},\qquad
 \mu=\max_{b\in B}\phi_B(b)>0.
\]
Give centres potential zero. Write $x_+=\max(x,0)$ and define
\[
 A=\sum_{t\in T}(\phi_T(t)-\phi_B(\pi(t)))_+,\qquad
 L=\sum_{t\in T}(\phi_B(\pi(t))-\phi_T(t))_+,\qquad
 R=\sum_{\substack{t\in T\\\pi(t)\in C}}\phi_T(t).
\]
Here $A$ counts aggregated withdrawals and $L$ aggregated compensation;
aggregation occurs before taking positive parts, so opposite regional
contributions on one path cancel.

\begin{theorem}\label{thm:hallmulti}
Under these hypotheses, $k\ge1$, $\max_T\phi_T=\mu$, and
\[
 A-L=\sum_{w\in W}\phi_B(\pi(w)),\qquad
 0\le A-L\le(k-1)\mu.
\]
If $q$ counts TH paths of positive compensation, then
\[
 q\ge\max\left\{0,\left\lceil\frac{A-(k-1)\mu}{\mu}\right\rceil\right\}
 \ge\max\left\{0,\left\lceil\frac{R-(k-1)\mu}{\mu}\right\rceil\right\}.
\]
Every such path is non-direct and ends in $B$. Consequently
\[
 f\ge k+\max\left\{0,\left\lceil\frac{A-(k-1)\mu}{\mu}\right\rceil\right\}.
\]
\end{theorem}
\begin{proof}
Fix a direct matching $m:B\to T$ saturating $B$. For every tight pair,
$m(D_j)\subseteq F_j$ and equal cardinality gives $m(D_j)=F_j$.
Thus $\phi_T(m(b))=\phi_B(b)$, and unmatched whites have potential
zero. The nonzero potential multisets on the two sides agree, giving
equal sums and maxima.

There cannot be $r$ HH centres with just $r-1$ HH sources, so $k\ge1$.
Exactly $r-k$ HH sources end at centres, leaving $k-1$ for the entire
exterior. Subtracting the HH contribution from the terminal bijection,
and using $x_+-(-x)_+=x$, gives the displayed conservation identity.
Each exterior HH endpoint contributes its summed potential once, at
most $\mu$, proving the upper bound. No source is allocated separately
to the regions containing its endpoint.

For every direct edge $tb$, membership $b\in D_j$ implies $t\in F_j$.
Therefore $\phi_T(t)\ge\phi_B(b)$, so positive compensation requires
a non-direct path. Its endpoint is exterior, and its amount is at most
$\mu$. Hence $q\mu\ge L\ge A-(k-1)\mu$. Every centre TH path contributes
its full source potential to $A$, giving $A\ge R$. The $k$ centre paths
and $q$ compensating paths are distinct non-direct paths, each using at
least one separate actual edge. Thus $f\ge k+q$.
\end{proof}

The normalization by $\mu$ is essential for overlapping regions. As an
endpoint countermodel, take
\[
 T=\{t_0,t_1,t_2,e\},\quad B=\{b_0,b_1,b_2\},\quad
 C=\{h_0,h_1\},\quad W=\{w\},
\]
with direct edges $t_ib_i$ for $i=0,1,2$, and bijection
\[
 t_0\mapsto h_1,\quad t_1\mapsto h_0,\quad t_2\mapsto b_2,
 \quad e\mapsto b_0,\quad w\mapsto b_1.
\]
Use $D_1=\{b_0,b_1\}$ and $D_2=\{b_1,b_2\}$, with their corresponding
white pairs $F_1,F_2$ and unit weights. The two regional triples
(withdrawals, TH compensation, HH compensation) are $(2,1,1)$ and
$(1,0,1)$. The same sole HH endpoint is counted twice. Subtracting
only $k-1=1$ from the unnormalized total would falsely require $1\ge2$.
The correct potentials give $\mu=2$, $A=3$, $L=1$ and
$A-L=(k-1)\mu=2$.

The centre-only version replaces $A$ by $R$. Since $R\le k\mu$, its
extra term is at most one, and equals one precisely when
$R>(k-1)\mu$. The full withdrawal bound can force more. For example,
take $T=\{t_1,t_2,t_3,e\}$, $B=\{b_1,b_2,b_3\}$, $C=\{h\}$,
$W=\varnothing$, direct edges $t_1b_1,t_2b_2,t_3b_3,t_2b_3$, and
$D=\{b_1,b_2\}$, $F=\{t_1,t_2\}$. The assignment
$t_1\mapsto h,t_2\mapsto b_3,t_3\mapsto b_1,e\mapsto b_2$ has
$k=\mu=R=1$ and $A=L=2$. The full bound gives $f\ge3$ for any path
realization satisfying the hypotheses, while the centre-only bound gives
$f\ge2$. Both examples concern endpoint deductions; no rectangular-grid
realization is asserted. Compatible physical routing and full completion,
the universal unit-weight $5/6$ improvement and the rectangular-grid
triple equality remain separate questions.


\paragraph{Verification.}
The results have the general proofs above. Independent finite controls
reconstruct both exchanges for every minimal set in all binary bipartite
adjacency matrices of total order at most eight, and in additional seeded
models. A separate complete subset enumeration checks all 256 interfaces
between two copies of $H$. Source and reproduction instructions are provided in the accompanying
\texttt{verification\_20261111.zip}; its README identifies the pinned
repository revision and complete verification reports. Additional controls check all 65,536 four-by-four witness interfaces
for complete cores, the paired formula for arbitrary seeded cores, and
the degree bounds by direct subset enumeration. Separate controls
check the general boundary predicates and abstract terminal bijections
for the simultaneous Hall budget. The new proofs
are not included in the earlier Lean formalization of the two-attachment
construction.

\paragraph{AI use.}
Haobo Ma and Wenlin Zhang used AI tools for search, double-checking and proofreading.

\begingroup
\small
\begin{thebibliography}{9}
\setlength{\itemsep}{2pt}
\bibitem{CF09} M. Chellali and O. Favaron,
On $k$-star-forming sets in graphs,
\emph{J. Combin. Math. Combin. Comput.} \textbf{68} (2009), 205--214.
\bibitem{Favaron12} O. Favaron,
Bounds on the Upper $k$-Domination Number and the Upper $k$-Star-Forming
Number of a Graph,
\emph{J. Combin. Math. Combin. Comput.} \textbf{80} (2012), 321--332.
\bibitem{Sahbi} R. Sahbi,
\emph{Upper $k$-Star-Forming Sets, $k$-Independence, and Upper Domination},
\href{https://arxiv.org/abs/2610.03785v2}{arXiv:2610.03785v2},
6 October 2026; subsequently withdrawn.
\bibitem{Bhattacharjee} D. Bhattacharjee,
\emph{Upper $k$-star-forming sets and $k$-independence in bipartite graphs},
version 1.1.0, 14 October 2026,
\href{https://doi.org/10.5281/zenodo.23360032}{doi:10.5281/zenodo.23360032}.
\bibitem{BPPP} F. Bock, J. Pardey, L. D. Penso and D. Rautenbach,
A bound on the dissociation number,
\emph{J. Graph Theory} \textbf{103} (2023), 661--673,
\href{https://doi.org/10.1002/jgt.22940}{doi:10.1002/jgt.22940}.
\bibitem{Li} P. Li,
\emph{A characterization of graphs attaining a lower bound on the
dissociation number},
\href{https://arxiv.org/abs/2610.35408v1}{arXiv:2610.35408v1}, 26 October 2026.
\end{thebibliography}
\endgroup
\end{document}
